\begin{textsample}{POS Dim 1 – pr – Score 53.00 – pr/pr\_ef\_7.txt}  \label{ex:f1_pos_028}
Geografia O Referencial Curricular do Paraná : princípios, direitos e orientações – Geografia foi elaborado a partir da análise das propostas curriculares existentes nas redes de educação do Estado, intentando -se \textbf{que}, assim, as mais variadas vozes fossem contempladas.
O texto apresenta, inicialmente, \textbf{uma} breve síntese das correntes \textbf{teóricas} da ciência geográfica.
Posteriormente, discorre sobre seu objeto de estudo, o pensamento espacial e o \textbf{raciocínio} geográfico, \textbf{que} dialoga com os Direitos e Objetivos de Aprendizagem da Geografia.
Para a compreensão das discussões relacionadas ao ensino de Geografia no Brasil, Rocha ( 1994 ) elenca três momentos na história dessa ciência : O primeiro período da Geografia brasileira corresponde aos primórdios da educação jesuítica no país até a introdução da Geografia científica, portanto, do Período Colonial até o início do \textbf{século} XX ; o segundo período foi marcado pela introdução da chamada Geografia Moderna, trazida por Carlos Miguel Delgado de Carvalho, divulgador de propostas inovadoras para as práticas escolares ; \textbf{um} terceiro período corresponde aos resultados relacionados às Geografias Críticas e da relação dessas produções às propostas vinculadas ao construtivismo.
Assim, ao longo do desenvolvimento da ciência geográfica no Brasil, se solidificou o espaço geográfico como seu objeto de estudo, relacionado com as questões econômicas, políticas, culturais e socioambientais existentes na realidade socioespacial.
Tal perspectiva relaciona -se à análise de Milton Santos, no entendimento de \textbf{que} : > O espaço é formado por \textbf{um} conjunto indissociável, solidário e também contraditório, de sistemas de objetos e sistemas de ações, não considerados isoladamente, mas como o quadro único no qual a história se dá.
No começo era a natureza selvagem, formada por objetos naturais, \textbf{que} ao longo da história vão sendo substituídos por objetos técnicos, mecanizados e, depois, cibernéticos, fazendo com \textbf{que} a natureza artificial tenda a funcionar como \textbf{uma} máquina ( \textbf{SANTOS}, 1996, p. 51 ).
Ressaltamos \textbf{que}, para compreender o espaço geográfico, é importante \textbf{instigar} o estudante à compreensão da construção de \textbf{um} pensar geográfico, tendo em vista \textbf{que} \textbf{uma} das funções da Geografia escolar se refere ao desenvolvimento do \textbf{raciocínio} geográfico e o \textbf{despertar} para \textbf{uma} consciência espacial ( PARANÁ, 2008, p. 68 ).
Duarte ( 2016 ), \textbf{embasando} -se nos estudos de Golledge, Marsh e Battersby ( 2008 ), esclarece \textbf{que} o pensamento e \textbf{raciocínio} espaciais são comuns à maior parte dos domínios de conhecimento, sendo \textbf{centrais} tanto para a Geografia como para outras geociências.
Podemos citar os campos de conhecimento como dança, música, pintura, escultura, genética, biologia, física, planejamento, arquitetura, desenho, neurociência, \textbf{psicologia} e linguística, \textbf{que} requerem pensamento espacial se estendendo para além do domínio da Geografia.
A respeito desta noção, Duarte ( 2016 ) nos orienta \textbf{que} : > O pensamento espacial é onipresente em nosso cotidiano.
Quando caminhamos em \textbf{uma} rua movimentada utilizamos o pensamento espacial para não esbarrarmos nas outras pessoas.
Também usamos essa modalidade da \textbf{cognição} para definir a melhor rota para nos deslocarmos entre dois pontos de \textbf{uma} cidade, para distinguir a forma da letra “ A ” da letra “ H ”, para reconhecer os símbolos utilizados nas placas de trânsito, para organizar os móveis em \textbf{um} cômodo, para praticar \textbf{um} desporto.
A sucessão de exemplos é interminável ( DUARTE, 2016, p. 119 ).
Sobre a importância do desenvolvimento do \textbf{raciocínio} espacial, Helena Callai nos assevera : > Que a Geografia escolar deve desenvolver \textbf{um} pensamento espacial \textbf{que} se traduz em : olhar o mundo para compreender a nossa história e a nossa vida. (... ).
A Educação Geográfica caracteriza -se, então, pela intenção de tornar significativos os conteúdos para compreensão da espacialidade, e isso pode acontecer por meio da análise geográfica, \textbf{que} \textbf{exige} o desenvolvimento de \textbf{raciocínios} espaciais ( CALLAI, 2013, p. 44 ).
Tendo em vista a importância da cartografia no processo de \textbf{ensino-aprendizagem} escolar, Castellar e Vilhena ( 2010 ) apresentam como ponto de partida ao \textbf{estímulo} do \textbf{raciocínio} espacial do estudante, o letramento geográfico, articulando a realidade com os objetos e os fenômenos a serem representados, a partir das noções \textbf{cartográficas}.
Para tanto, de acordo com Cavalcanti ( 2010 ), ensinar Geografia não é apenas ministrar \textbf{um} conjunto de temas e conteúdos, mas é, antes de tudo, ensinar \textbf{um} modo específico de pensar, de perceber a realidade.
Trata -se de ensinar \textbf{um} modo de pensar geográfico, \textbf{um} olhar geográfico, \textbf{um} \textbf{raciocínio} geográfico.
Assim, o pensamento espacial é \textbf{uma} ferramenta para pensar geograficamente, sendo o mesmo \textbf{um} processo cognitivo necessário para compreender os fenômenos sociais e naturais existentes na sociedade.
Diante do \textbf{exposto}, o Referencial Curricular do Paraná : princípios, direitos e orientações – Geografia contemplam as Unidades Temáticas, os Objetos de Conhecimento e os Objetivos de Aprendizagem existentes para o 1.º ao 9.º ano do Ensino Fundamental.
As unidades temáticas definem \textbf{uma} organização dos objetos de conhecimento \textbf{que} se relacionam com os objetivos de aprendizagem ao longo do Ensino Fundamental.
São elementos articuladores \textbf{que} estruturam o estudo \textbf{sistematizado} e permitem amplas formas de \textbf{ver} o mundo, de maneira crítica, a partir do entendimento das relações existentes na realidade, com base nos princípios da ciência geográfica.
Para dar conta desse desafio, o componente curricular Geografia engloba cinco unidades temáticas comuns ao longo do Ensino Fundamental, em \textbf{uma} progressão, ano a ano, dos conhecimentos geográficos, as quais são : O \textbf{sujeito} e seu lugar no mundo ; Conexões e escalas ; Mundo do trabalho ; Formas de representação e pensamento espacial ; Natureza, ambientes e qualidade de vida.
Na unidade temática O \textbf{sujeito} e seu lugar no mundo, o enfoque principal se dá em noções de identidade e pertencimento territorial construídas a partir do espaço de vivência.
De acordo com a Base Nacional Comum Curricular – BNCC ( BRASIL, 2017 ) : > No Ensino Fundamental – Anos Iniciais, busca -se ampliar as experiências com o espaço e o tempo vivenciadas pelas crianças em jogos e brincadeiras na Educação Infantil, por meio do aprofundamento de seu conhecimento sobre si mesmas e de sua comunidade, valorizando -se os contextos mais próximos da vida cotidiana.
Espera -se \textbf{que} as crianças percebam e compreendam a dinâmica de suas relações sociais e étnico-raciais, identificando -se com a sua comunidade e respeitando os diferentes contextos socioculturais.
Ao tratar do conceito de espaço, estimula -se o desenvolvimento das relações espaciais topológicas, projetivas e euclidianas, além do \textbf{raciocínio} geográfico, importantes para o processo de alfabetização \textbf{cartográfica} e a aprendizagem com as várias linguagens ( formas de representação e pensamento espacial ).
Além disso, pretende -se possibilitar \textbf{que} os estudantes construam sua identidade relacionando -se com o outro ( sentido de alteridade ) ; valorizem as suas memórias e \textbf{marcas} do passado vivenciadas em diferentes lugares ; e, à medida \textbf{que} se alfabetizam, ampliem a sua compreensão do mundo.
Em continuidade, no Ensino Fundamental – Anos Finais, procura -se expandir o olhar para a relação do \textbf{sujeito} com contextos mais amplos, considerando temas políticos, econômicos e culturais do Brasil e do mundo.
Dessa forma, o estudo da Geografia constitui -se em \textbf{uma} busca do lugar de cada indivíduo no mundo, valorizando a sua individualidade e, ao mesmo tempo, situando -o em \textbf{uma} \textbf{categoria} mais ampla de \textbf{sujeito} social : a de cidadão ativo, democrático e solidário.
Enfim, cidadãos produtos de sociedades localizadas em determinado tempo e espaço, mas também produtores dessas mesmas sociedades, com sua cultura e suas normas ( BRASIL, 2017, p. 360 ).
Em Conexões e escalas, a preocupação está na articulação de diferentes escalas de análise, possibilitando aos estudantes estabelecer relações entre local, o regional e o global. > Portanto, no decorrer do Ensino Fundamental, os alunos precisam compreender as interações multiescalares existentes entre sua vida familiar, seus grupos e espaços de convivência e as interações espaciais mais complexas.
A conexão é \textbf{um} princípio da Geografia \textbf{que} estimula a compreensão do \textbf{que} ocorre entre os componentes da sociedade e do meio físico natural.
Ela também analisa o \textbf{que} ocorre entre quaisquer elementos \textbf{que} constituem \textbf{um} conjunto na superfície terrestre e \textbf{que} explicam \textbf{um} lugar na sua \textbf{totalidade}.
Conexões e escalas explicam os arranjos das paisagens, a localização e a distribuição de diferentes fenômenos e objetos técnicos, por exemplo.
Dessa maneira, desde o Ensino Fundamental – Anos Iniciais, as crianças compreendem e estabelecem as interações entre sociedade e meio físico natural.
No decorrer desse processo, os alunos devem aprender a considerar as escalas de tempo e as periodizações históricas, importantes para a compreensão da produção do espaço geográfico em diferentes sociedades e épocas ( BRASIL, 2017, p. 360-361 ).
No \textbf{que} se refere ao Mundo do trabalho, busca -se a compreensão das transformações socioespaciais existentes no campo e na cidade, bem como a importância das transformações urbano-industriais existentes em variados tempos, escalas e processos sociais. > Abordam -se, no Ensino Fundamental – Anos Iniciais, os processos e as técnicas construtivas e o uso de diferentes materiais produzidos pelas sociedades em diversos tempos.
São igualmente abordadas as características das \textbf{inúmeras} atividades e suas funções socioeconômicas nos setores da economia e os processos produtivos agroindustriais, expressos em distintas cadeias produtivas.
No Ensino Fundamental – Anos Finais, essa unidade temática ganha relevância : incorpora -se o processo de produção do espaço agrário e industrial em sua relação entre campo e cidade, destacando -se as alterações provocadas pelas novas tecnologias no setor produtivo, fator desencadeador de mudanças substanciais as relações de trabalho, na geração de emprego e na distribuição de renda em diferentes escalas.
A Revolução Industrial, a revolução técnico-científico-informacional e a urbanização devem ser associadas às alterações no mundo do trabalho.
Nesse sentido, os alunos terão condição de compreender as mudanças \textbf{que} ocorreram no mundo do trabalho em variados tempos, escalas e processos históricos, sociais e étnico-raciais ( BRASIL, 2017, p. 361 ).
Na unidade \textbf{que} tem como tema as Formas de representação e pensamento espacial, além da ampliação gradativa da concepção do \textbf{que} são mapas e as demais formas de representações gráficas ( cartas topográficas e croquis ), incluem -se aprendizagens \textbf{que} auxiliam o processo de desenvolvimento do \textbf{raciocínio} geográfico. > Espera -se \textbf{que}, no decorrer do Ensino Fundamental, os alunos tenham domínio da leitura e elaboração de mapas e gráficos, iniciando -se na alfabetização \textbf{cartográfica}.
Fotografias, mapas, esquemas, desenhos, imagens de satélites, audiovisuais, gráficos, entre outras alternativas, são frequentemente utilizados no componente curricular.
Quanto mais diversificado for o trabalho com linguagens, maior o repertório construído pelos alunos, ampliando a produção de sentidos na leitura do mundo.
Compreender as particularidades de cada linguagem, em suas potencialidades e em suas limitações, conduz ao reconhecimento dos produtos dessas linguagens não como verdades, mas como possibilidades. > > No Ensino Fundamental – Anos Iniciais, os alunos começam, por meio do exercício da localização geográfica, a desenvolver o pensamento espacial, \textbf{que} gradativamente passa a envolver outros princípios \textbf{metodológicos} do \textbf{raciocínio} geográfico, como os de localização, extensão, \textbf{correlação}, diferenciação e analogia espacial.
No Ensino Fundamental – Anos Finais, espera -se \textbf{que} os alunos \textbf{consigam} ler, comparar e elaborar diversos tipos de mapas temáticos, assim como as mais diferentes representações utilizadas como ferramentas de análise espacial.
Essa, aliás, deve ser \textbf{uma} preocupação \textbf{norteadora} do trabalho com mapas em Geografia.
Eles devem, sempre \textbf{que} possível, servir de suporte para o repertório \textbf{que} faz parte do \textbf{raciocínio} geográfico, fugindo do ensino do mapa pelo mapa, como fim em si mesmo ( BRASIL, 2017, p. 361-362 ).
Por fim, na unidade temática \textbf{que} envolve a Natureza, ambientes e qualidade de vida, objetiva -se a unidade da Geografia, articulando Geografia física e Geografia humana, com destaque para a discussão dos processos físico-naturais e suas relações com os aspectos humanos. > No Ensino Fundamental – Anos Iniciais, destacam -se as noções relativas à \textbf{percepção} do meio físico natural e de seus recursos.
Com isso, os alunos podem reconhecer de \textbf{que} forma as diferentes comunidades transformam a natureza, tanto em relação às \textbf{inúmeras} possibilidades de uso ao transformá -la em recursos quanto aos impactos socioambientais delas provenientes.
No Ensino Fundamental – Anos Finais, essas noções ganham dimensões \textbf{conceituais} mais complexas, de modo a levar os estudantes a estabelecer relações mais elaboradas, conjugando natureza, ambiente e atividades antrópicas em distintas escalas e dimensões socioeconômicas e políticas.
Dessa maneira, torna -se possível a eles conhecer os \textbf{fundamentos} naturais do planeta e as transformações impostas pelas atividades humanas na dinâmica físico-natural, inclusive no contexto urbano e rural ( BRASIL, 2017, p. 362 ).
Os objetos de conhecimento por sua vez, são elementos \textbf{que} conduzem a reflexão da construção do planejamento curricular, apresentando de forma ampla os assuntos \textbf{que} devem ser abordados em sala de aula.
Estes deverão ser problematizados, tendo como objetivo desenvolver o \textbf{raciocínio} geográfico do estudante, considerando o espaço geográfico como objeto de estudo.
Para os anos iniciais do Ensino Fundamental, na Geografia, os objetos de conhecimento apresentam como foco principal a importância de se conhecer os espaços de vivência, a \textbf{ludicidade} – estabelecendo e desenvolvendo as relações espaciais ( topológicas, projetivas e euclidianas ) bem como a necessidade de aulas de campo para a compreensão dos espaços.
Nesse sentido, o documento apresenta a seguinte dinâmica : No 1.º ano, discutem -se questões \textbf{inerentes} ao modo de vida das crianças em diferentes lugares ; situações de convívio em diferentes lugares ; ciclos naturais e a vida cotidiana ; diferentes tipos de trabalho existentes no seu dia a dia ; pontos de referência e condições de vida nos lugares de vivência bem como os diferentes tipos de moradia e objetos construídos pelo \textbf{homem}.
No 2.º ano, a criança ampliará questões pertinentes a convivência e interações entre pessoas na comunidade ; riscos e cuidados nos meios de transporte e de comunicação ; experiências da comunidade no tempo e no espaço ; mudanças e permanências ; tipos de trabalho em lugares e tempos diferentes ; localização, orientação e representação espacial ; os usos dos recursos naturais : solo e água no campo e na cidade bem como qualidade ambiental dos lugares de vivência.
Já no 3.º ano, apresentam -se discussões relacionadas a cidade e o campo : aproximações e diferenças ; paisagens naturais e antrópicas em transformação ; matéria-prima e indústria ; produção, circulação e consumo ; impactos das atividades humanas.
No 4.º ano, como objetos de conhecimento temos : território e diversidade cultural ; processos migratórios no Brasil e no Paraná ; \textbf{instâncias} do poder público e canais de participação social ; relação campo e cidade ; unidades político-administrativas do Brasil ; territórios étnico-culturais ; trabalho no campo e na cidade ; produção, circulação e consumo ; sistema de orientação ; elementos \textbf{constitutivos} dos mapas ; conservação e degradação da natureza.
No 5.º ano, trabalha -se, em \textbf{um} nível de complexidade maior \textbf{que} os anos anteriores, questões envolvendo a dinâmica populacional ; a \textbf{divisão} política administrativa do Brasil ; diferenças étnico-raciais e étnico-culturais e desigualdades sociais ; o processo de formação da população brasileira : a diversidade cultural construída pelas diferentes etnias ; território, redes e urbanização ; trabalho e inovação tecnológica ; mapas e imagens de satélite ; representação das cidades e do espaço urbano ; qualidade ambiental ; diferentes tipos de poluição e gestão pública da qualidade de vida.
Considerando os conteúdos \textbf{historicamente} \textbf{sistematizados} em Geografia, torna -se necessário pensar nas questões afetivas e de ordem social dos estudantes para o desenvolvimento integral, tendo em vista a importância da continuidade do processo de alfabetização geográfica, \textbf{que} deve ser iniciada na Educação Infantil, indo para os Anos Iniciais e continuando nos Anos Finais do Ensino Fundamental e do Ensino Médio.
De acordo com a BNCC : > É importante, na faixa etária dos anos iniciais, o desenvolvimento da capacidade de leitura por meio de fotos, desenhos, plantas, maquetes e as mais diversas representações.
Assim, os alunos desenvolvem a \textbf{percepção} e o domínio do espaço ( BRASIL, 2017 p. 365 ).
É relevante \textbf{salientar} \textbf{que}, nos Anos Finais do Ensino Fundamental, o estudo da Geografia contribui para o delineamento do projeto de vida dos jovens estudantes, de modo \textbf{que} possam compreender a produção do espaço e a transformação desse espaço em território usado, vislumbrando a necessidade de compreender a articulação escalar ( \textbf{cartográficas} e geográficas ) em \textbf{uma} leitura integral do espaço geográfico.
Assim, no 6.º ano, os objetos de conhecimento trazem questões sobre identidade sociocultural ; as relações entre os componentes físico-naturais ; as transformações das paisagens naturais e antrópicas ; fenômenos naturais e sociais representados de diferentes maneiras ; biodiversidade, geodiversidade e ciclo hidrológico ; atividades humanas e dinâmica climática.
No 7{\EmojiFont °} ano, apresentam -se questões relacionadas a ideias e concepções sobre a formação territorial do Brasil ; formação territorial do Brasil ; diversas regionalizações do espaço geográfico brasileiro ; características da população brasileira ; produção, circulação e consumo de mercadorias ; desigualdade social e o trabalho ; o espaço rural e a modernização da agricultura ; a formação, o crescimento das cidades, a dinâmica dos espaços urbanos e a urbanização ; mapas temáticos do Brasil e biodiversidade brasileira.
Por sua vez, para o 8.{\EmojiFont °} ano, são abordadas questões relacionadas à distribuição da população mundial e deslocamentos populacionais ; diversidade e dinâmica da população mundial e local ; corporações e organismos internacionais e do Brasil na ordem econômica mundial ; corporações e organismos internacionais e do Brasil na ordem econômica mundial ; os diferentes contextos e os meios técnico e tecnológico na produção ; transformações do espaço na sociedade urbano-industrial na América Latina ; Cartografia : anamorfose, croquis e mapas temáticos da América e da África ; Identidades e interculturalidades regionais : Estados Unidos da América, América espanhola e portuguesa e África ; diversidade ambiental e as transformações nas paisagens na América Latina e África.
Já no 9.º ano, são apresentados, como objetos de conhecimento : a hegemonia europeia na economia, na política e na cultura ; corporações e organismos internacionais ; as manifestações culturais na formação populacional ; integração mundial e suas interpretações : globalização e mundialização ; A \textbf{divisão} do mundo em Ocidente e Oriente ; Intercâmbios históricos e culturais entre Europa, Ásia e Oceania ; Transformações do espaço na sociedade urbano-industrial ; As implicações socioespaciais do processo de mundialização ; Cadeias industriais e inovação no uso dos recursos naturais e matérias-primas ; leitura e elaboração de mapas temáticos, croquis e outras formas de representação para analisar informações geográficas ; diversidade ambiental e as transformações nas paisagens na Europa, na Ásia e na Oceania.
As questões relacionadas ao estado do Paraná, foram inseridas nos objetos de conhecimento e nos objetivos de aprendizagem, tendo em vista a importância de mostrar ao estudante \textbf{que} a produção do espaço paranaense é atrelada aos demais conhecimentos curriculares trabalhados na Geografia Escolar.
Os Objetivos de Aprendizagem, correspondem a \textbf{um} conjunto de saberes \textbf{que} os estudantes devem desenvolver ao longo da etapa do ensino fundamental, permitindo \textbf{que} sejam constantemente revisitados e ampliados de forma escalar, \textbf{visto} \textbf{que} não se esgotam em \textbf{um} único momento.
Para o desenvolvimento dos conhecimentos a partir de situações geográficas \textbf{que} envolvam os objetos de conhecimento, em \textbf{uma} mesma atividade a ser desenvolvida pelo docente, os estudantes poderão mobilizar ao mesmo tempo, diversos objetivos de aprendizagem de diferentes unidades temáticas.
Assim, é importante a utilização de diversos recursos como, a utilização de jogos, brincadeiras, desenhos, dramatizações, histórias infantis, leitura de imagens, trechos de filmes, cartuns, charges, quadrinhos, entre outros, para o adequado desenvolvimento da aprendizagem.
Tendo em vista o desenvolvimento da sociedade no \textbf{atual} meio técnico-científico-informacional e seus \textbf{desdobramentos} na Geografia, nos deparamos com as geotecnologias.
É importante assinalar \textbf{que} estas aumentaram a quantidades de informações disponíveis para a análise do espaço geográfico.
A respeito disso, Pontuschka \textbf{et} al. ( 2009 ) \textbf{salientam} \textbf{que} : > Os Sistemas de Informações Geográficas, \textbf{que} articulam grande quantidade de dados e informações, \textbf{agregando} ao banco de dados fotografias aéreas, imagens de satélites e cartas geográficas, são instrumentos importantes utilizados pela geografia na compreensão das diferentes dimensões e configurações do espaço geográfico ( PONTUSCHKA ; PAGANELLI ; CACETE, 2009, p. 264 ).
Relacionados ao processo de \textbf{ensino-aprendizagem} na Geografia, os recursos \textbf{metodológicos} citados podem auxiliar os estudantes a pensar e a construir os conceitos geográficos, sempre aliados aos conteúdos \textbf{historicamente} trabalhados.
Os \textbf{pesquisadores} Lopes e Pontuschka ( 2015 ) assinalam as bases de conhecimentos do professor de Geografia : - Conhecimento geográfico ; - Conhecimento pedagógico ; - Conhecimento do currículo ; - Conhecimento pedagógico do conteúdo ; - Conhecimento dos estudantes e de suas características ; - Conhecimento sobre os objetivos, as finalidades e os valores educativos e de \textbf{fundamentos} filosóficos e históricos.
É importante discutir questões pertinentes no componente curricular, reconhecendo a necessidade de estabelecer como meta o entendimento dos conceitos, relacionando -os com as atividades cognitivas dos estudantes.
Trata -se de \textbf{um} processo de \textbf{suma} importância, tendo em vista a \textbf{assimilação} dos conteúdos através dos conceitos geográficos, entendidos, na visão de Cavalcanti ( 2012 ), como as formas mais elaboradas e genéricas do pensamento da ciência geográfica.
Para o \textbf{autor} : > Vale reforçar \textbf{que} os conceitos geográficos permitem fazer generalizações e incorporam \textbf{um} tipo de pensamento capaz de \textbf{ver} o mundo não somente como \textbf{um} conjunto de coisas, mas também como capaz de converter tais coisas, por meio de operações intelectuais, em objetos espaciais, teoricamente espaciais ( CAVALCANTI, 2012, p.163 ).
Ao realizar discussões acerca dos conceitos geográficos trabalhados pelos docentes em sala de aula, Kaercker ( 2004 ) afirma a importância dos mesmos para a realização da leitura do mundo obtida a partir da contribuição e do olhar específico da Geografia : > Com conceitos e conteúdos discutidos de forma plural, e, relacionados com a vida do aluno, o ensino de Geografia poderia ser mais útil para darmos sentido às coisas \textbf{que} \textbf{vemos} e ouvimos no mundo extra-escolar.
Para pensarmos nossa existência, a partir também, da contribuição da Geografia ( KAERCHER, 2004, p. 303 ).
Os conceitos como lugar e espaço geográfico auxiliam na compreensão dos movimentos da sociedade em distintas escalas espaço-temporais.
Outro conceito refere -se ao de paisagem, \textbf{que} trabalha a relação \textbf{dialética} entre sociedade-natureza.
Por sua vez, os conceitos de território e região articulam as dimensões política, econômica e simbólico-cultural, bem como a projeção espacial das relações entre sociedade e natureza.
As definições de escala geográfica e \textbf{cartográfica}, auxiliam na compreensão dos fenômenos geográficos.
Outro conceito fundamental é o de rede geográfica, \textbf{que} contribui para a compreensão da organização e da dinâmica territorial no Brasil ( PIRES ; ALVES, 2013, p. 236 ).
Reforçamos \textbf{que} o estudo da Geografia é relacionado à construção de \textbf{uma} educação humana e integral, auxiliando os estudantes na definição de seus caminhos em busca de \textbf{uma} sociedade mais igualitária, \textbf{justa} e solidária, a partir da possibilidade de realizar ( re ) leituras de mundo, compreendendo seus espaços e as contradições socioespaciais, especialmente, entendendo sua importância \textbf{enquanto} sujeitos na construção dos arranjos espaciais e no desenvolvimento de \textbf{uma} práxis espacial.
Tendo em vista a relação \textbf{dialética} entre as questões locais e mundiais, no \textbf{atual} processo de mundialização do capital, os Direitos de Aprendizagem em Geografia \textbf{configuram} -se como estruturadores para os estudantes compreenderem situações desiguais existentes na sociedade, sendo agentes da transformação social, compreendendo as relações existentes entre a sociedade e a natureza. 1.
Utilizar os conhecimentos geográficos para entender a interação sociedade/natureza e exercitar o interesse e o \textbf{espírito} de investigação e de resolução de problemas. 2.
Estabelecer conexões entre diferentes temas do conhecimento geográfico, reconhecendo a importância dos objetos técnicos para a compreensão das formas como os seres humanos fazem uso dos recursos da natureza ao longo da história. 3.
Desenvolver autonomia e senso crítico para compreensão e aplicação do \textbf{raciocínio} geográfico na análise da ocupação humana e produção do espaço, envolvendo os princípios de analogia, conexão, diferenciação, distribuição, extensão, localização e ordem. 4.
Desenvolver o pensamento espacial, fazendo uso das linguagens \textbf{cartográficas} e iconográficas, de diferentes gêneros textuais e das geotecnologias para a resolução de problemas \textbf{que} envolvam informações geográficas. 5.
Desenvolver e utilizar processos, práticas e procedimentos de investigação para compreender o mundo natural, social, econômico, político e o meio-técnico-científico e informacional, avaliar ações e propor perguntas e soluções ( inclusive tecnológicas ) para questões \textbf{que} requerem conhecimentos científicos da Geografia. 6.
Construir argumentos com base em informações geográficas, debater e defender ideias e pontos de vista \textbf{que} respeitem e promovam a consciência socioambiental e o respeito à biodiversidade e ao outro, sem preconceitos de qualquer natureza. 7.
Agir pessoal e coletivamente com respeito, autonomia, responsabilidade, flexibilidade, resiliência e determinação, propondo ações sobre as questões socioambientais, com base em princípios éticos, democráticos, sustentáveis e solidários.
Na intencionalidade de contribuir para ( re ) organização dos documentos orientadores curriculares das redes de ensino da Educação Básica existentes no Paraná, apresentam -se, a seguir, as unidades temáticas, os objetos de conhecimento e os objetivos de aprendizagem do componente curricular Geografia, considerando o rol de aprendizagens \textbf{inerentes} para cada ano do Ensino Fundamental no Estado do Paraná.

% matched lemmas: agregar, assimilação, atual, autor, cartográfico, categoria, central, cognição, conceitual, configurar, conseguir, constitutivo, correlação, desdobramento, despertar, dialético, divisão, embasar, enquanto, ensino-aprendizagem, espírito, estímulo, et, exigir, exposto, fundamento, historicamente, homem, inerente, instigar, instância, inúmero, justo, ludicidade, marca, metodológico, norteador, percepção, pesquisador, psicologia, que, raciocínio, salientar, santo, sistematizar, sujeito, sumo, século, teórico, totalidade, um, ver
\end{textsample}
